\documentclass[11pt]{article}

\usepackage[utf8]{inputenc}
\usepackage[T1]{fontenc}
\usepackage{amsmath,amssymb}
\usepackage{booktabs}
\usepackage{enumitem}
\usepackage{hyperref}
\usepackage[margin=1in]{geometry}

\title{\textbf{Context Management and Language Memoization:\\
Dynamic Programming for Recursive Language-Model Agents}}

\author{Adarsh Vatsa}
\date{Draft -- April 2026}

\begin{document}

\maketitle

\begin{abstract}
Long-context language-model systems are often framed as search or retrieval
problems. A more precise framing is that the agent is solving a context
management problem under a fixed token budget. Recursive Language Models
(RLMs) formalize one way to let a model inspect, decompose, and recursively
call itself over snippets of an oversized context. However, this does not
exhaust the design space: a well-built agent runtime with token-budget
feedback, addressable context, isolated subagents, compaction, and
self-reprompting can implement the same general strategy. This note argues
that the more important missing layer is not recursion itself, but reusable
recursive execution: a dynamic-programming-style memoization system for
language-model subproblems. Instead of merely caching final answers, such a
system should store structured intermediate work, source spans, dependency
trails, confidence metadata, and reuse modes that allow future agents to
replay, verify, compose, narrow, or rule out prior work. We also describe the
first concrete implementation in this repository: a small structural
memoization layer supporting exact replay, negative memoization, compositional
coverage over context intervals, and JSON persistence.
\end{abstract}

\section{Introduction}

Long-context agents are increasingly asked to answer questions over inputs
that exceed the context window of the underlying language model. One response
is to build longer-context models. Another is to treat the full context as
external to any single model call and require the agent to decide what subset
of context should be inspected at each step.

The second view makes the underlying systems problem clear. The agent's job is
not simply to ``search'' the corpus. Its job is to manage scarce attention.
It must decide when to delegate, when to compact, when to inspect a region,
when to spawn a new instance of itself, and when the evidence it has collected
is sufficient to answer.

This paper-style note develops that argument and then motivates a stronger
cache abstraction for recursive agents. The proposed abstraction is not
ordinary semantic caching of the form
\[
    \text{query} + \text{context} \rightarrow \text{answer}.
\]
It is closer to dynamic programming over language-model subproblems: once a
recursive subproblem has been solved over a context region, future calls should
be able to reuse that work exactly, partially, compositionally, or as search
guidance.

\section{Recursive Context Management}

Consider an agent given a question $Q$ and a context $C$ that is many times
larger than its usable context window. The general solution is an iterative
context-management loop:

\begin{enumerate}[leftmargin=*]
    \item Detect that $C$ exceeds the usable prompt budget.
    \item Choose a context-management move: delegate slices to subagents,
    compact state, inspect a selected region, or ask a new instance of itself
    with selected context and state.
    \item Receive compact, structured outputs.
    \item Decide whether the evidence is sufficient to answer $Q$.
    \item If not, repeat with a better-targeted prompt or context selection.
\end{enumerate}

This loop does not require one particular scaffold. It requires an
orchestrator with feedback and control: approximate token counts, addressable
context, isolated subagents, structured returns, and the ability to construct
new prompts for future model instances.

\section{RLM as a Formalization}

Zhang, Kraska, and Khattab's Recursive Language Models (RLMs) formalize one
version of this loop. Their abstraction treats long prompts as part of an
external environment and lets the model programmatically examine, decompose,
and recursively call itself over snippets of the prompt~\cite{zhang2026rlm}.

This is a useful formal substrate, but it is not a fundamentally different
class of operation from a well-built long-context agent runtime. A capable
runtime with token-budget feedback, addressable context, isolated subagents,
compaction, and self-reprompting can implement the same general strategy.

The distinction can be summarized as follows:

\begin{itemize}[leftmargin=*]
    \item RLM is a disciplined substrate for recursive context management.
    \item A strong agent runtime can be functionally equivalent if it exposes
    the same control over context selection, delegation, and compaction.
    \item The deeper unsolved problem is not merely whether an agent can
    recurse over context, but whether the system can avoid re-solving work it
    has already done.
\end{itemize}

\section{Search as Candidate Generation}

A model attending over a bounded chunk can often outperform classical search
inside that chunk. It can resolve paraphrase, interpret local evidence, and
synthesize answers in ways BM25 or vector similarity cannot.

This does not imply that search should be discarded. Classical search and
indexing remain valuable as cheap candidate generators and context-addressing
mechanisms. Their purpose is not to reason better than the model. Their purpose
is to cheaply propose where model attention should be spent.

Useful low-cost signals include token counts, chunk identifiers, source spans,
file or page metadata, headings, lexical matches, vector candidates, cached
summaries, and prior subagent results.

The target architecture is therefore index-assisted recursive attention:
cheap indexing proposes candidate regions; bounded model calls read those
regions; the orchestrator reasons over compact returns.

\section{From Semantic Cache to Language Memoization}

The current repository implements a useful semantic cache. It stores answers
for query-context pairs, scopes reuse by source or dataset identity, verifies
semantic equivalence before replay, and persists cache state across runs.

However, ordinary answer caching is too coarse for recursive agents. A future
recursive call may not need the same final answer. It may need a supporting
fact, a local summary, a proof fragment, a pointer to a likely region, a
record that some region was already ruled out, or a partial aggregate over a
subset of the context.

The stronger abstraction is a memoized language artifact:

\begin{center}
\begin{tabular}{ll}
\toprule
\textbf{Field} & \textbf{Meaning} \\
\midrule
task & Subproblem or question solved \\
context\_scope & Document, chunk, span, or set of spans \\
result & Answer, fact, summary, decision, or pointer \\
evidence & Source spans supporting the result \\
dependencies & Subcalls used to produce the result \\
confidence & Confidence score or verifier status \\
reusable\_as & Replay, fact, hint, aggregate, or ruled-out region \\
\bottomrule
\end{tabular}
\end{center}

This shifts the system from binary cache hit/miss behavior to a reuse planner.

\section{Reuse Modes}

A memoized recursive system should support several reuse modes:

\begin{description}[leftmargin=1.25in,style=nextline]
    \item[Replay.] The same task has already been solved over the same context
    scope, so the prior result can be returned directly.

    \item[Verify.] A prior result likely applies, but a cheap verifier or
    deterministic check should confirm that reuse is safe.

    \item[Compose.] A larger task can reuse solved subregions and only call
    models for missing or ambiguous regions.

    \item[Narrow.] Prior work identifies where the answer likely lives.

    \item[Rule out.] Prior work identifies regions where the answer likely
    does not live.

    \item[Aggregate.] Prior partial answers can be combined into a larger
    answer.
\end{description}

The model's role changes accordingly. Instead of attending over all context
from scratch, model calls become semantic operators for memoized work:

\begin{itemize}[leftmargin=*]
    \item Are two tasks equivalent?
    \item Does an old result entail what the new query needs?
    \item Does an old context span overlap enough to reuse?
    \item Can partial results compose into the larger answer?
    \item Is the prior result trustworthy enough to skip a model call?
\end{itemize}

\section{A DP-Shaped Solver}

The target solver has the structure of dynamic programming, but with semantic
matching and verification in place of exact symbolic equality:

\begin{verbatim}
Solve(question Q, context C):
  if memo has an exact solved entry for Q over C:
      return it

  if memo has useful partial entries for Q over regions of C:
      reuse those entries
      identify missing or ambiguous regions
  else:
      identify candidate regions with indexes, summaries, or decomposition

  for each missing region:
      call a bounded subagent with Q and that region
      store the subagent result as a memo entry

  aggregate partial results
  verify answer and evidence
  store Solve(Q, C)
  return answer
\end{verbatim}

This is different from ordinary semantic caching because the goal is not only
answer replay. The goal is work avoidance and search guidance throughout a
recursive reasoning process.

\section{Relationship to the Current Repository}

The existing implementation already contains several pieces needed for this
direction:

\begin{itemize}[leftmargin=*]
    \item \texttt{source\_chunk\_hash} protects direct query reuse over exact
    source contexts.
    \item \texttt{data\_scope\_hash} protects retrieval-cache reuse across
    different document sets.
    \item The LLM Sniper verifies semantic equivalence before replaying a
    cached answer.
    \item Knowledge extraction begins moving from answer blobs toward
    fact-level reuse.
    \item Persistence allows work to survive across runs.
\end{itemize}

The architectural shift is:
\[
    \text{cache answers} \quad \longrightarrow \quad
    \text{memoize recursive subproblem work}.
\]

This requires first-class intermediate artifacts, evidence spans, dependency
trails, verifier metadata, and explicit reuse modes.

\section{Concrete V0 Implementation}

The repository now includes a first implementation of this idea in
\texttt{language\_memoization.py}. This module is deliberately smaller than the
full semantic-cache system. It does not call an LLM, compute embeddings, or try
to solve fuzzy semantic equivalence. Instead, it implements the deterministic
substrate that a semantic layer can build on. The main
\texttt{SemanticCacheController} now owns a \texttt{MemoStore}, persists it with
the existing cache state, and exposes a controller-level
\texttt{memoized\_subproblem()} method for recursive workflows.

The v0 goal is:

\begin{quote}
Given a task and an addressable context interval, determine which parts of the
work can be replayed, which parts have already been ruled out, which hints are
available, and which intervals still require model calls.
\end{quote}

\subsection{Core Types}

\paragraph{TaskSpec.}
\texttt{TaskSpec} describes the subproblem being solved. It contains the
natural-language prompt, a task type, an output contract, and optional
constraints. V0 task matching is exact after normalization: whitespace is
collapsed and text is lowercased. This is intentionally conservative.

\paragraph{ContextScope.}
\texttt{ContextScope} is a half-open interval:
\[
    [\texttt{start}, \texttt{end})
\]
inside a corpus and document. The unit can be \texttt{chunk}, \texttt{char},
\texttt{token}, or any caller-defined coordinate system. A scope also carries
an optional \texttt{content\_hash}; if two scopes have different non-empty
hashes, they are treated as incompatible. This prevents memoized work from one
version of a document being reused against another.

\paragraph{MemoEntry.}
\texttt{MemoEntry} stores one reusable unit of model work. It contains the task,
scope, result, evidence spans, dependencies, confidence, verifier status, and
one or more reuse labels. The initial reuse labels are:

\begin{center}
\begin{tabular}{ll}
\toprule
\textbf{Label} & \textbf{Meaning} \\
\midrule
\texttt{exact\_answer} & Can directly answer the same task/scope \\
\texttt{supporting\_fact} & Can support a future answer \\
\texttt{search\_hint} & Points future work toward a likely region \\
\texttt{aggregation\_component} & Can compose into a larger aggregate \\
\texttt{ruled\_out} & This region was already checked and failed \\
\bottomrule
\end{tabular}
\end{center}

\paragraph{MemoStore.}
\texttt{MemoStore} is the in-memory table. It supports adding answer entries,
adding negative entries, planning reuse, and saving/loading JSON state.

\subsection{Planner Semantics}

The central method is:

\begin{verbatim}
plan = store.plan_reuse(task, scope)
\end{verbatim}

The planner follows four deterministic rules:

\begin{enumerate}[leftmargin=*]
    \item If an \texttt{exact\_answer} exists for the same normalized task and
    exact same compatible scope, return an exact replay plan.
    \item Otherwise, collect all entries with the same normalized task and a
    compatible overlapping scope.
    \item Split them into reusable answer/aggregate entries, negative
    \texttt{ruled\_out} entries, and \texttt{search\_hint} entries.
    \item Treat answer, aggregate, supporting-fact, and ruled-out entries as
    coverage. Subtract their intervals from the requested scope to produce
    \texttt{missing\_scopes}.
\end{enumerate}

This gives the orchestrator a concrete action plan. If
\texttt{missing\_scopes} is empty, the requested scope is fully covered by
prior work. If not, the orchestrator only needs to spawn subagents for the
missing intervals.

\subsection{Controller Integration}

The controller integration has two pieces.

First, \texttt{SemanticCacheController.memoized\_subproblem()} lets callers
solve scoped subproblems through the memo planner. The caller provides a
\texttt{solver(task, scope)} function. In production, that solver can call a
subagent or model. In tests, it is a deterministic function. The controller
calls the solver only for \texttt{missing\_scopes}, stores each new result as a
\texttt{MemoEntry}, and then stores a composed exact entry for the larger scope
when coverage is complete.

Second, \texttt{search()} is memo-first. It mirrors legacy final-answer cache
rows and legacy knowledge triples into memo entries, then checks the memo graph
for exact scoped replay before falling back to document retrieval and answer
synthesis. The older data structures remain loadable, but they are compatibility
inputs rather than separate reuse systems.

The structural flow is:

\begin{verbatim}
plan = memo_store.plan_reuse(task, requested_scope)

if plan.has_exact_replay:
    return prior_answer

for scope in plan.missing_scopes:
    result = solver(task, scope)
    memo_store.add_answer(...) or memo_store.add_negative(...)

if all intervals are covered:
    memo_store.add_answer(task, requested_scope, composed_answer)
\end{verbatim}

\subsection{DuckDB Backend}

The default \texttt{MemoStore} is an in-memory dictionary with JSON
persistence. That is appropriate for tests and small local runs, but it scans
entries linearly. The repository therefore includes
\texttt{DuckDBMemoStore}, an optional local backend for larger benchmark runs.

DuckDB is used narrowly: it stores durable structured memo rows, raw context
chunks, dependency edges, and performs fast candidate lookup. Python still
performs interval subtraction and reuse planning. This keeps the database role
simple and keeps planner behavior identical across memory and DuckDB backends.
The model is not given unrestricted SQL execution. Instead, the controller
exposes structured operations such as lexical chunk search and range fetch, then
passes bounded candidate packets to the model.

The important table columns are:

\begin{center}
\begin{tabular}{ll}
\toprule
\textbf{Column} & \textbf{Purpose} \\
\midrule
\texttt{entry\_id} & Stable memo entry identity \\
\texttt{task\_signature} & Normalized task lookup key \\
\texttt{corpus\_id}, \texttt{document\_id} & Corpus and document isolation \\
\texttt{scope\_start}, \texttt{scope\_end} & Interval lookup \\
\texttt{unit} & Coordinate system, e.g. chunk or char \\
\texttt{content\_hash} & Document-version guard \\
\texttt{reusable\_as\_json} & Reuse labels \\
\texttt{entry\_json} & Full round-trippable memo entry \\
\texttt{context\_chunks} & Raw chunk text keyed by corpus, document, and index \\
\bottomrule
\end{tabular}
\end{center}

The key lookup is:

\begin{verbatim}
SELECT entry_json
FROM memo_entries
WHERE task_signature = ?
  AND corpus_id = ?
  AND document_id = ?
  AND unit = ?
  AND content_hash = ?
  AND scope_start < requested_end
  AND scope_end > requested_start
\end{verbatim}

This returns only memo entries for the same task and overlapping context
interval. The planner then classifies them as exact replay, reusable coverage,
negative coverage, or hints.

Usage is intentionally parallel to the in-memory store:

\begin{verbatim}
from language_memoization import DuckDBMemoStore

store = DuckDBMemoStore("memo.duckdb")
store.add_answer(task, scope, "$12.4M")
plan = store.plan_reuse(task, larger_scope)
store.close()
\end{verbatim}

When \texttt{solve\_with\_memo()} receives chunked context and the memo store is
DuckDB-backed, the controller also persists those chunks in a
\texttt{context\_chunks} table. The practical retrieval surface is:

\begin{verbatim}
store.fetch_context_range(
    corpus_id="my-corpus",
    document_id="doc-1",
    start=0,
    end=4,
)

store.search_context_chunks(
    "launch codeword",
    corpus_id="my-corpus",
    document_id="doc-1",
)

controller.context_candidate_packets("launch codeword", scope=scope)
\end{verbatim}

This gives the orchestrator an addressable context substrate without turning
retrieval into raw model-authored SQL. A future planner can ask for structured
operations such as ``fetch this range'' or ``search chunks for these terms'',
while the implementation remains parameterized and auditable.

This backend is useful when memo entries become numerous enough that JSON
round-trips and linear scans are no longer comfortable, but a full server
database such as Postgres would be operationally excessive.

\subsection{Example}

Suppose an agent must find a financial metric over chunks 0 through 10. The
memo store already knows that chunks 0 through 2 were checked and did not
contain the metric, and that chunks 5 through 7 contain a useful partial
answer. The planner returns only the uncovered intervals:

\begin{verbatim}
from language_memoization import (
    ContextScope, MemoStore, REUSE_AGGREGATION_COMPONENT, TaskSpec
)

task = TaskSpec(
    prompt="Find the Q1 ARR",
    task_type="financial_metric_lookup",
    output_contract="short factual answer",
)

def scope(start, end):
    return ContextScope(
        corpus_id="finance",
        document_id="model-a",
        start=start,
        end=end,
        unit="chunk",
        content_hash="document-version-1",
    )

store = MemoStore()
store.add_negative(task, scope(0, 2), reason="metric not found")
store.add_answer(
    task,
    scope(5, 7),
    "$12.4M appears in the KPI table",
    reusable_as=(REUSE_AGGREGATION_COMPONENT,),
)

plan = store.plan_reuse(task, scope(0, 10))

# plan.missing_scopes is:
#   [ContextScope(start=2, end=5), ContextScope(start=7, end=10)]
\end{verbatim}

The important point is that the system does not ask ``do we have the final
answer?'' It asks ``which parts of this recursive computation are already
solved?''

\subsection{Negative Memoization}

Negative entries are first-class. If a subagent has already inspected a region
and found nothing, that result is useful. It prevents future agents from
spending attention on the same failed region. In search-heavy recursive tasks,
this can be as valuable as caching positive facts.

For example, a \texttt{ruled\_out} entry over chunks 0--2 counts as solved
coverage for the same task. A later query over chunks 0--10 only needs to
inspect the rest.

\subsection{Current Clean Architecture}

The current implementation is now memo-first. The older semantic-cache and
knowledge mechanisms still exist for backward compatibility, but they no longer
own independent reuse decisions in the main \texttt{search()} path. They are
mirrored into memo entries:

\begin{itemize}[leftmargin=*]
    \item legacy final answers become \texttt{exact\_answer} memo entries;
    \item legacy knowledge triples become \texttt{supporting\_fact} memo entries;
    \item FAISS, lexical memo search, and DuckDB chunk search are candidate
    generators;
    \item exact replay, partial composition, semantic verifier reuse, evidence,
    and dependency edges are represented through the memo graph.
\end{itemize}

The practical boundary is \texttt{reuse\_candidate\_packets()}. It exposes the
candidate generators as bounded packets while keeping the memo graph as the
single source of truth for reuse.

\begin{verbatim}
packet = controller.reuse_candidate_packets(query, scope=scope)

packet["memo_entries"]     # reusable work
packet["context_chunks"]   # raw evidence candidates
\end{verbatim}

Q3 is used when language judgment is required: solving scoped tasks,
aggregating partial memo entries, or verifying semantic reuse. It is not the
cache, and it does not write raw SQL.

\section{Research Positioning}

The clean distinction from RLM is not that recursive context management is
new. It is that recursive context management benefits from persistent, scoped,
semantic memoization.

\begin{center}
\begin{tabular}{ll}
\toprule
\textbf{Layer} & \textbf{Role} \\
\midrule
RLM & Recursively navigate and process huge context \\
Memoization & Remember and reuse recursive subproblem work \\
Provenance & Determine which subanswers are trustworthy \\
Planning & Use prior subanswers to guide future decomposition \\
\bottomrule
\end{tabular}
\end{center}

Put differently:

\begin{quote}
RLM gives recursive execution. Language memoization gives reusable recursive
execution.
\end{quote}

\section{Practical Next Steps}

The current repository now has a practical MVP of the automatic path:

\begin{verbatim}
answer = controller.solve_with_memo(
    query,
    chunks,
    solver=local_q3_solver,
    document_id="synthetic-doc",
    chunk_size=2,
    aggregate_fn=local_q3_aggregator,
)
\end{verbatim}

This method decomposes a document into chunk windows, checks the memo store for
each scoped subproblem, calls the supplied solver only for missing windows,
passes reusable partial answers to an aggregator, stores the full-scope answer,
and replays that full answer on later calls. It also accepts a
\texttt{reuse\_verifier} callback. That verifier may approve reuse from a
different prompt over the same scope, producing a new exact memo entry whose
dependency points back to the source memo entry. This is the first safe step
toward semantic dynamic programming: fuzzy reuse is possible, but only when a
verifier explicitly allows it.

The live smoke test exercises this path with DuckDB and a local MLX Qwen model:

\begin{verbatim}
python scripts/live_dp_memo_smoke.py \
  --reset-db \
  --backend duckdb \
  --duckdb-path benchmark_artifacts/live_dp_memo.duckdb
\end{verbatim}

The benchmark scripts deliberately avoid small repository-level generation caps
on the normal path. The \texttt{--max-tokens} and
\texttt{--aggregate-max-tokens} flags remain available for explicit stress
tests, but omitted flags leave generation length to the MLX backend or to
\texttt{RLMS\_MLX\_MAX\_TOKENS} if the operator sets it.

The expected shape is:

\begin{itemize}[leftmargin=*]
    \item first solve: two scoped model calls and one aggregation call;
    \item exact replay: zero scoped model calls and zero aggregation calls;
    \item verifier-approved semantic replay: zero solver and aggregation calls
    for a different prompt over the same scope;
    \item persistence: reopening the DuckDB file preserves exact replay.
\end{itemize}

The same DP memoization path also has a small NoLiMa fixture benchmark:

\begin{verbatim}
python scripts/run_dp_memo_nolima.py \
  --reset-db \
  --max-cases 2 \
  --max-haystacks 1 \
  --max-samples 4 \
  --lengths 250 \
  --depth-intervals 2 \
  --chunk-words 60 \
  --chunk-size 2
\end{verbatim}

This benchmark uses the repository's local NoLiMa needle and haystack fixtures.
It writes predictions, bridge rows, a manifest, and a DuckDB memo store under
\url{benchmark_artifacts/dp_memo_nolima/}. In the current smoke run, the
four fixture samples cover two needles across two depths, and every sample is
immediately replayed from DuckDB after the first solve.
\texttt{--max-cases} and \texttt{--max-haystacks} define the source slice;
\texttt{--max-samples} is only an early stop inside that generated slice.

A more direct benchmark for memoized reuse is the shared-context fixture:

\begin{verbatim}
python scripts/run_dp_memo_shared_context.py \
  --reset-db \
  --sentences-per-chunk 2 \
  --chunk-size 1
\end{verbatim}

This benchmark asks several related questions over one document. The system
first solves a reusable fact-extraction task over the document chunks and stores
those facts in DuckDB. Each question is then answered by giving the local model
compact memo candidate packets and asking it to select the useful entries. This
tests the complementary path directly: cheap memo retrieval proposes candidates;
the model reads structured memo content and chooses the answer. The current run
answers all five fixture questions correctly and verifies zero-call replay for
the memoized fact layer.

The DuckDB backend now also materializes the memo graph explicitly.
Dependencies are mirrored into a \texttt{memo\_edges} table with parent ID,
child ID, edge type, timestamp, and metadata. The initial read APIs are:

\begin{itemize}[leftmargin=*]
    \item \texttt{graph\_edges()}
    \item \texttt{children(entry\_id)}
    \item \texttt{parents(entry\_id)}
    \item \texttt{lineage(entry\_id)}
\end{itemize}

This keeps the existing deterministic replay path unchanged while making the
dependency graph auditable and queryable.

\subsection{Current Instrumentation}

The memo substrate now exposes telemetry at two levels.

First, every \texttt{MemoReusePlan} can report JSON-safe coverage telemetry:

\begin{itemize}[leftmargin=*]
    \item requested scope length;
    \item covered length;
    \item missing length;
    \item coverage ratio;
    \item covered intervals;
    \item missing scopes;
    \item exact, reusable, negative, and hint entry IDs;
    \item fragment-kind counts for all candidates and covering candidates;
    \item evidence-span and dependency-edge counts.
\end{itemize}

This matters because the benchmark should not only report whether the final
answer was correct. It should report how much work the memo graph avoided. A
typical row can now distinguish:

\begin{verbatim}
initial coverage: 0.00
final coverage:   1.00
reused windows:   2 / 3
coverage kinds:   supporting_fact=1, ruled_out_region=2
\end{verbatim}

Second, the store exposes \texttt{stats()}:

\begin{itemize}[leftmargin=*]
    \item total memo entries;
    \item counts by fragment kind;
    \item counts by result type;
    \item counts by verifier status;
    \item counts by reuse mode;
    \item dependency edge count;
    \item evidence span count;
    \item average confidence.
\end{itemize}

The fragment kind is planner-facing. Instead of asking a model to infer the
role of a memo row from loose text, the packet tells it whether the row is an
\texttt{exact\_answer}, \texttt{supporting\_fact},
\texttt{aggregation\_component}, \texttt{search\_hint}, or
\texttt{ruled\_out\_region}. This is a small schema choice, but it is central
to the ``split cache'' framing: the value of the system depends on storing the
right fragments, not just storing more answer strings.
Because reuse-plan telemetry includes fragment-kind counts, an ablation can
now distinguish coverage from exact answer replay, supporting facts,
aggregation components, and negative memoization instead of reporting only one
opaque coverage number.

Candidate packet text can still be bounded for prompt construction, but
truncation is no longer silent. Memo packets include
\texttt{result\_chars}, \texttt{result\_truncated},
\texttt{evidence\_count}, and per-evidence \texttt{text\_truncated} fields.
Raw context packets include \texttt{text\_chars} and
\texttt{text\_truncated}. Passing \texttt{None} for the relevant max-character
argument disables truncation. This avoids a common failure mode in agentic
systems: hidden context budgets silently removing the evidence the planner
needed.

\subsection{Long-Context Benchmark Path}

The project now includes a NoLiMa workload inspector:

\begin{verbatim}
python scripts/inspect_long_context_workload.py \
  --needle-set-path benchmark_data/nolima/needlesets/needle_set.json \
  --haystack-dir benchmark_data/nolima/haystack/rand_shuffle_long \
  --lengths 32K,64K,128K,256K \
  --depth-intervals 4 \
  --max-cases 4 \
  --max-haystacks 1 \
  --chunk-words 1000 \
  --chunk-size 2
\end{verbatim}

This script does not call a model. It estimates how many samples, chunks,
solver windows, and cold model calls a benchmark slice implies. For the example
slice above, the inspector reports 64 samples and 3,904 estimated cold model
calls. The 64K, 128K, and 256K settings force recursive decomposition rather
than direct prompting.

This is the correct benchmark direction for this repo because NoLiMa
\cite{modarressi2025nolima} minimizes
literal overlap between the question and the needle. A vector or lexical
retriever can still serve as a candidate generator, but it should not be able
to solve the task by string matching alone. The memo layer should be evaluated
on whether repeated domain-scoped calls progressively fill a reusable fragment
graph that reduces future work.

The first official-data smoke surfaced a fragment-design issue. Asking every
scoped chunk to produce the final answer caused over-filtering: the relevant
chunk contained a local fact such as ``Katie lives next to the Kiasma museum,''
while the question asked which character had been to Helsinki. A direct-answer
subcall returned \texttt{NOT\_FOUND}. The better fragment was not a final
answer, but a reusable evidence fragment. The runner now supports:

\begin{verbatim}
--solver-mode evidence
\end{verbatim}

In this mode, scoped calls extract broad named-entity and relationship facts,
and the aggregator performs the final inference. On the same official 8K
NoLiMa smoke sample, this changed the answer from an incorrect \texttt{NoLiMa}
to the correct answer, \texttt{Katie}, while preserving the same memo behavior:
21 scoped calls, one aggregate call, full final coverage, and zero-call exact
replay.

A broader 8K slice later exposed a sharper two-hop failure. A question asked
which character had been to \texttt{Uusimaa}; the relevant slice said
\texttt{Katie lives next to the Kiasma museum}. The first evidence extractor
still over-filtered and returned \texttt{NOT\_FOUND}, so the aggregator never
received the bridge fact. The runner now makes evidence mode stricter about
recall: if a slice contains any named person, place, landmark, title, code, or
relationship, the extractor should emit a fact instead of \texttt{NOT\_FOUND}.
On the targeted failed sample this changed the answer back to
\texttt{Katie}. The tradeoff is clear: recall-oriented fragments improve
two-hop correctness but produce more supporting facts and higher latency.

The warm-start run over the same DuckDB store then answered the same official
sample with zero scoped model calls and zero aggregation calls. This gives the
minimal real benchmark lifecycle: cold execution builds the fragment graph;
warm execution replays from it.

The next test used overlapping length sweeps. By opting into stable corpus,
document, and source-hash identities, a 16K solve can become usable prefix work
for a later 32K solve over the same source. In the 32K run, the planner started
with 81 of 161 chunks already covered, skipped 40 windows, and called the model
only for the expanded tail. The answer remained correct. This is the first
practical partial-composition result: the larger task reused a solved subregion
rather than replaying an identical full answer.

The same pattern held as the official slice grew:

\begin{center}
\begin{tabular}{lrrrr}
\toprule
Run & Initial coverage & Model calls & Aggregate calls & Correct \\
\midrule
16K cold & 0.000 & 41 & 1 & yes \\
32K partial & 0.503 & 41 & 1 & yes \\
64K partial & 0.502 & 81 & 1 & yes \\
128K partial & 0.501 & 161 & 1 & yes \\
128K warm & 1.000 & 0 & 0 & yes \\
256K partial & 0.579 & 234 & 1 & yes \\
256K warm & 1.000 & 0 & 0 & yes \\
\bottomrule
\end{tabular}
\end{center}

The 128K partial run reused 321 of 641 chunks from the existing memo graph,
processed only the 320 missing chunk positions, and then stored a full-scope
answer. A fresh-process warm replay over the same DuckDB file answered with
zero solver calls, zero aggregation calls, and about 150 ms of wall-clock time.
The 256K run extended the same pattern: it started with 641 of 1108 chunks
covered, skipped 320 windows, solved 234 missing windows, answered correctly,
and then warm-replayed in about 135 ms with zero model calls.
An attempted 512K probe over the same single-haystack slice produced the same
1108 chunk scope and exact-replayed immediately. That is a useful benchmark
limitation: once the selected haystack source is exhausted, increasing the
requested length label no longer creates a harder workload. Larger experiments
therefore need either larger source files or multiple haystacks/cases, not just
a larger nominal length flag.
This does not prove general agent quality, but it does show the intended
substrate behavior: as a corpus grows in overlapping stages, prior scoped work
remains usable.

\subsection{What the Benchmark Does and Does Not Prove}

The NoLiMa path is useful because it stress-tests a real failure mode: long
contexts where literal overlap is weak and naive chunk answering can miss
indirect evidence. It is not, by itself, a complete measure of real-world agent
performance. The benchmark has a particular shape: a synthetic long context, a
hidden needle, a short answer, no evolving workspace, and no naturally repeated
user workflow unless the experiment creates one.

This distinction matters for memoized RLMs. The central claim is not that the
system should overfit a single long-context QA format. The claim is that
repeated, domain-bounded knowledge work creates reusable intermediate
subproblems:

\begin{itemize}[leftmargin=*]
    \item the same corpus is queried repeatedly;
    \item questions overlap without being identical;
    \item useful evidence fragments survive across tasks;
    \item old negative searches can prevent repeated dead ends;
    \item partial work over a smaller scope can become useful inside a larger
    scope.
\end{itemize}

A cold, single-shot benchmark mostly measures first-pass search. That can
understate the value of a persistent memo graph, because the graph becomes more
valuable as a workload continues. Conversely, real deployments introduce risks
that synthetic benchmarks may understate: stale fragments, changed documents,
permission boundaries, ambiguous user goals, and unsafe semantic reuse.

The right evaluation target is therefore a workload trace, not only a single
question-answer pair. A good benchmark for this substrate should report:

\begin{itemize}[leftmargin=*]
    \item final accuracy;
    \item model calls, wall-clock time, and token cost;
    \item cold-start performance;
    \item warm replay performance;
    \item partial-overlap reuse across growing or changing corpora;
    \item stale-memory failure rate after document updates;
    \item which fragment kinds were actually reused.
\end{itemize}

In this framing, NoLiMa is one instrument in the evaluation suite. It helps
surface fragment-design problems. The stronger evidence for this project comes
from workload-style runs where the memo graph is filled once, then reused across
related tasks or larger overlapping scopes.

\subsection{Staleness and Scope Invalidation}

Persistent memoization introduces a safety problem that ordinary one-shot
benchmarks rarely test: old work can become stale. Content hashes prevent many
accidental reuses, but a practical agent also needs an explicit operational
control for document updates. The memo store now exposes:

\begin{verbatim}
store.invalidate_scope(scope, reason="document updated")
\end{verbatim}

This rejects all non-rejected memo entries that overlap the supplied scope. The
entries are not deleted. They remain available for audit and lineage, but exact
replay, partial composition, and planner candidate generation ignore them.
Invalidation also propagates through dependency parents by default: if a
derived answer depends on a changed fragment, that derived answer is rejected
even when its own scope does not directly overlap the edited chunk. This matters
because a long-lived memo graph should behave like a versioned working memory,
not an append-only bag of facts that are always trusted.

Passing a scope with an empty content hash invalidates compatible entries for
that document range regardless of version. Passing a specific content hash only
invalidates entries compatible with that version. This gives callers a direct
way to handle document edits, corpus refreshes, or user corrections without
destroying the historical record.

There are two useful versioning modes:

\begin{itemize}[leftmargin=*]
    \item Use strict content hashes when any document change should force a
    clean recomputation.
    \item Use stable or empty content hashes plus explicit
    \texttt{invalidate\_scope()} when the system should reuse unaffected
    windows after a localized update.
\end{itemize}

The second mode is closer to a real mutable workspace. A regression test now
covers this path: after chunk 1 changes, the old full answer and old chunk-1
fragment are rejected, chunks 0, 2, and 3 remain reusable, and the next solve
calls the solver only for chunk 1.

The repository also includes a deterministic mutable-workspace benchmark:

\begin{verbatim}
python scripts/run_dp_memo_mutable_workload.py --reset-db
\end{verbatim}

The trace solves a four-chunk operational runbook, updates one chunk,
invalidates that scope, solves the updated version, and then reopens DuckDB for
warm replay. The current run has the intended shape:

\begin{verbatim}
v1_model_calls: 4
invalidated_entries: 2
v2_model_calls: 1
v2_initial_coverage_ratio: 0.75
warm_model_calls: 0
\end{verbatim}

This is a tiny benchmark, but it tests a qualitatively important workload
pattern that single-shot long-context QA does not: persistent memory must remain
useful across localized document updates.

\subsection{Composition Fix}

One implementation bug surfaced while using the substrate on the long-context
path. The document-level \texttt{solve\_with\_memo()} loop used to skip a
window only when that exact window had an exact replay entry. That missed a
central dynamic-programming case: a larger memo fragment may already fully
cover a smaller window.

The loop now skips any window whose reuse plan is complete. This allows a
window to be avoided if it is covered by:

\begin{itemize}[leftmargin=*]
    \item an exact window answer;
    \item a larger solved region;
    \item supporting facts;
    \item an aggregation component;
    \item a ruled-out region.
\end{itemize}

This change preserves the old exact replay path but makes partial composition
real in the automatic document solver.

Concrete next steps for the repository are:

\begin{enumerate}[leftmargin=*]
    \item Run longer official NoLiMa slices with Q3 and inspect coverage,
    reused-window ratio, answer accuracy, and wall-clock cost.
    \item Extend \texttt{solve\_with\_memo()} from fixed chunk windows to
    structure-aware windows based on sections, records, pages, or functions.
    \item Add verifier calls only when reuse would skip meaningful model work.
    \item Add semantic task matching on top of the structural task signature.
    \item Build ablations comparing no memoization, exact answer replay,
    semantic answer replay, fact reuse, and full memoized subproblem reuse.
\end{enumerate}

\section{Conclusion}

Recursive context management is a powerful way to make fixed-context language
models operate over inputs much larger than their context windows. But
recursion alone does not prevent repeated work. The more durable systems
contribution is a language memoization graph: a persistent, scoped, verified
record of recursive subproblems and their reusable intermediate results.

Such a system would let an orchestrator solve a long-context task once and
then reuse the resulting subproblem graph whenever later tasks overlap with
that work.

\begin{thebibliography}{9}

\bibitem{zhang2026rlm}
Alex L. Zhang, Tim Kraska, and Omar Khattab.
\newblock Recursive Language Models.
\newblock \emph{arXiv preprint arXiv:2512.24601}, 2026.
\newblock \url{https://arxiv.org/abs/2512.24601}.

\bibitem{modarressi2025nolima}
Ali Modarressi, Hanieh Deilamsalehy, Franck Dernoncourt, Trung Bui,
Ryan A. Rossi, Seunghyun Yoon, and Hinrich Schuetze.
\newblock NoLiMa: Long-Context Evaluation Beyond Literal Matching.
\newblock \emph{International Conference on Machine Learning}, 2025.
\newblock \url{https://arxiv.org/abs/2502.05167}.

\bibitem{bai2025longbenchv2}
Yushi Bai, Shangqing Tu, Jiajie Zhang, Hao Peng, Xiaozhi Wang, Xin Lv,
Shulin Cao, Jiazheng Xu, Lei Hou, Yuxiao Dong, and Juanzi Li.
\newblock LongBench v2: Towards Deeper Understanding and Reasoning on
Realistic Long-context Multitasks.
\newblock \emph{ACL}, 2025.
\newblock \url{https://longbench2.github.io/}.

\bibitem{zhang2024infinitebench}
Xinrong Zhang, Yingfa Chen, Shengding Hu, Zihang Xu, Junhao Chen, Moo Hao,
Xu Han, Zhen Thai, Shuo Wang, Zhiyuan Liu, and Maosong Sun.
\newblock InfiniteBench: Extending Long Context Evaluation Beyond 100K Tokens.
\newblock \emph{ACL}, 2024.
\newblock \url{https://github.com/OpenBMB/InfiniteBench}.

\end{thebibliography}

\end{document}
